\begin{papercard}{halterman2025}
\cf{Stages} \stagetag{L} (a staged evaluation of whether an LLM follows a real codebook, with label-free tests before hand-labelling), \stagetag{C} (a semi-structured codebook format usable by humans and models: definition, clarification, negative clarification, examples). Political Analysis article (online 2025); arXiv version 2 used. Added from the gap search of 23 October 2026 as a pilot structure for D9.

\cf{Question} Whether off-the-shelf LLMs follow real codebook operationalisations \pc{halterman2025}{1}{there is limited empirical evidence on whether “off-the-shelf” LLMs faithfully follow real-world codebook operationalizations}, tested on three codebooks with up to 142 labels and long definitions.

\cf{Five stages} Stage 0, codebook preparation in a semi-structured format for humans and machines \pc{halterman2025}{3}{We suggest a semi-structured format with standardized components including definitions, clarifications, and examples that can be used by both humans and machines}; Stage 1, label-free behavioural tests before any hand-labelling \pc{halterman2025}{3}{Before investing in hand-labeling, we suggest researchers use our proposed label-free behavioral tests}; Stage 2, zero-shot evaluation on a small hand-coded set; Stage 3, error analysis with ablations and manual reading \pc{halterman2025}{4}{We suggest thorough error analysis using ablation experiments, behavioral tests requiring labels, and manual analysis of outputs}; Stage 4, supervised fine-tuning if zero-shot is inadequate.

\cf{Behavioural tests} Examples: does the model return only labels defined in the codebook \pc{halterman2025}{12}{Does the LLM only return labels defined in the codebook?}; can it recall a definition; are predictions stable when the order of categories changes; does it apply exclusion criteria consistently when document and codebook are modified \pc{halterman2025}{12}{Across the four combinations of (modi- fied, original) × (document, codebook), are the LLM predictions consistently cor- rect?}.

\cf{Findings} Open-weight 7--12B models complied poorly with codebooks zero-shot; instruction tuning on human-coded examples improved performance by up to 55\,\% \pc{halterman2025}{3}{we show that supervised instruction tuning—updating LLM weights directly from human-coded examples—can substantially improve performance (by up to 55\%)}.

\cf{Limits (card)} Small open models only; multi-class single-label document classification, not span annotation; political-science codebooks.

\cf{Use for this thesis} (1) A ready pilot structure for D9: write the codebook once with definition, clarification, negative clarification and example per label (the methodology already plans example and counter-example); run label-free tests first (only valid labels, definition recall, order stability, exclusion rules such as ``qualification alone does not establish minimisation''); then a small zero-shot evaluation; then error analysis. (2) The negative-clarification format fits the stance boundaries (qualified vs.\ minimised; reassurance vs.\ measure). (3) Fine-tuning is the fallback the thesis cannot afford, which is a further argument for keeping the model as a screener.
\end{papercard}
